\documentclass[12pt]{article}
\usepackage{graphicx} % Required for inserting images
\usepackage[a4paper, top=2.4cm, bottom=2.4cm, left=2.4cm, right=2.4cm]{geometry}
\usepackage{times} % Times New Roman
\setlength{\parindent}{0pt}
\usepackage[spanish]{babel}
\usepackage{url}

\begin{document}

\begin{center}
    {\bfseries\uppercase{Ciberseguridad, seguridad Informática y Seguridad de la Información}} \\
\textit{D. F. Velásquez Pichilla}\\
\textit{7690-16-3882 Universidad Mariano Gálvez}\\
\textit{Seminario de Tecnología de Información}\\
\textit{dvelasquezp4@miumg.edu.gt}\\
\end{center}

\textbf{URL REPOSITORIO LA TEX: } \\
https://github.com/DiegoVelasq98/Ciberseguridad-seguridad-informatica.git

\vspace{1em}
\noindent\textbf{Resumen}\\
La seguridad en general es un tema tan importante en la actualidad debido al avance que ha tenido la teconologia, un claro ejemplo es el poder acceder y hacer todo lo que se realizaba de forma presencial a un enfoque digitial, desde cualquier dispositivo con acceso a internet.
De tal forma no termina el proceso alli, ya que las empresas deben de garantizar a los usuarios que dicha informacion es segura y no podra ser vulnerada ni comprometida, ya que esto afectaria a ambas partes, desde la preferencia del servicio, y la otra en posibles perdidas.
Es por eso que siempre se debe tener un buen control de seguridad a nivel empresa, que cubra desde lo fisico, hasta el software , desde que la comunicacion con el cliente empieza, hasta con el fin de dicha informacion para tratar en el sistema, sin olvidar al personal , ya que al aplicar en conjunto todo eso, se obtiene el resultado de los sistemas de informacion, la cual es , como su nombre lo indica, la informacion. \\

\noindent\textbf{Palabras clave: } ciberseguridad, seguridad, disponibilidad, ciberataque

\vspace{1em}
\noindent\textbf{Ciberseguridad} \\
La ciberseguridad es aquella que se refiere a todas las practicas, normativas, gobernanzas, etc, que ayuden a prevenir ciberataques, ya que hoy en dia la informacion es la parte esencial de cada organizacion, la cual siempre se debe tener el maximo de cuidado y seguridad. \\ \\
Ya que en esta era moderna donde manda todo lo digital , la dependencia de la informacion y la conectividad, es importante resguardar desde equipos, aplicaciones de software, redes, activos criticos y todo aquello que pueda ser vulnerado o amenazado.\\\\
Ya que una de las obligaciones de las empresas digitales es garantizar a sus usuarios la proteccion total de sus datos alojados en su sistema ya sea por control interno de sus politicas y lo que dictamina la ley acerce de los datos. Durante las cuales deben tener rapidas respuestas cuando se llegue a presentar un ciberataque, ya que se debe contar con planes de prevencion y de correcion para mitigar los riesgos al minimo. Ya que nadie esta a salvo de no ser atacado, sumado a que los ciberataques cada vez son mas frecuentes, costosos y que obligan a las organizaciones a ser proactivos y a dedicar mas presupuesto a esta area vital.
\\\\
El tener medidas de cibereguridad en una organizacion dan como resultado:
\begin{itemize}
    \item Reduccion del costo de brechas
    \item Conformidad normativa
    \item Mitigar ciberamenazas en el desarrollo
\end{itemize}


\textbf{Seguridad Informatica} \\
Son todas aquellas medidas, normativas o gobernanzas que se realizan para proteger a los sistemas informaticos en una organizacion y el resultado de dichos sistemas, es decir , la informacion que se maneja para las operaciones de una organizacion. Uno de sus principales objetivos es que se cumpla el aseguramiento de la informacion, evitando de aquellos accesos que no esten autorizados, perdida de la misma o interrupciones de procesos que puedan afectar a la misma. \\\\
Esto abarca desde el enfoque fisico , contemplando desde equipos, dispositivos, servidores, luego pasando por el interno, es decir el software, contemplando a los sistemas operativos, programas, etc, para finalmente a los medios de conectividad, es decir la red interna de la empresa, todo esto garantizando que no tengan accesos no deseados, permisos, roles y todo aquello que implique el manejo del sistema.\\\\
Dandonos a entender que la seguridad no solo va enfocada a sistemas, si no de todo aquello que hace posible su funcionamiento, desde el personal, redes, el mismo sistema, equipos, componentes, todo aquello que pueda ser vulnerable o amenazado. Y de aca es que se llegan a tener alternativas como lo pueden ser firewalls, antivirus, controladores de trafico, analizadores de dispositivos en red,reglamentos, backups, etc.\\\\
Algunas de las amenazas mas comunes podremos mencionar:
\begin{itemize}
    \item No contar con backups
    \item Actualizaciones pendientes
    \item Acceso fisico no autorizado
    \item Legacy sistems
    \item Fallas de autenticacion multifactor internos
    \item  No tener segmentada la red
\end{itemize}

\textbf{Seguridad de la Informacion } \\
Es aquella que busca que se permita el correcto acceso y gestion de la informacion, basandose en tres principios fundamentales, los cuales son, disponibilidad, integridad y confidencialidad, los cuales se denominan la triada de la informacion, las cuales aplicarse correctamente , brindan la seguridad en los datos.\\
Siendo que cada una tiene como fin:
\begin{itemize}
    \item\textbf{Confidencialidad: }{No acceder a datos que no se esten autorizados}
    \item\textbf{Integridad: }{Que la informacion sea completa y precisa}
    \item\textbf{Disponibilidad: }{Que se pueda acceder a la misma cuando se necesite}
\end{itemize}
Cada uno de estos es importante ya que son una especie de normativas, reglamentos, sugerencias, gobernanzas o todo aquello que dictamine como debe ser resguardada y manipulada la informacion para evitar que sea vulnerable o expuesta a amenazas. \\ 

Siendo que se hace enfasis en que no va solo enfocada en lo informatico, sino en todo aquello que lo compone y  da como resultado la informacion, es decir, va desde, personas, equipos, red, hardware, software, medidas de seguridad fisicas, medidas de seguridad logicas, etc. Haciendo referencia que para aplicar cada una de estas medidas se deben de contar con apoyo, es decir las gobernanzas que ya rigen las normas, que garanticen como deben ser abordadas, para poder evitar que ocurran problemas, como lo pueden ser desde el acceso solo de personal autorizado a ciertas areas, prohibir el ingreso de alimentos en el area de trabajo, bloquear dhcp en la red, roles de acceso, etc. \\


\noindent\textbf{Observaciones y comentarios} 
\begin{itemize}
    \item{Es importante que las empresas siempre piensen en la seguridad de los datos que manejan, ya que nadie esta excepto de ataques o fallas que puedan comprometer la informacion, ya que es prefereble invertir en esto desde el principio antes de que ocurran las cosas a cuando pasan y el gasto sea mayor que el inicial}
    \item{Es importante que las empresas tengan medidas que no vayan solo desde sistemas informaticos, sino que es importantes que estas tengan contemplando todo aquello que hace posible su funcionamiento, como personal, red, equipos, etc}
\end{itemize}

\noindent\textbf{Conclusiones} 
\begin{itemize}
    \item{Se llego a la conclusion de que los ataques cada vez son mas innvoadores, lo cual obliga a las empresas a no solo aplicar medidas de seguridad, sino el darles tambien seguimiento a dichas medidas}
    \item {Se llego a la conclusion de que el tener buena seguridad en todos los niveles ayuda a que sea suceptible a vulnerabilidades o amenazas ya que en ocasiones si dicha informacion es comprometida, puede afectar a una empresa desde pausar operaciones hasta  quiza su cese.}
    \item{Se llego a la conclusion de que la seguridad es responsabilidad de todos dentro de una organizacion, ya que no sirve de mucho contar con buenos sistemas y herramientas de proteccion si las personas no conocen los riesgos o no hacen un uso adecuado de los mismos.}

\end{itemize}


\begin{thebibliography}{9}

\bibitem{Microsoft2025}
Microsoft. (s.f.). \textit{What is information security (InfoSec)}. Recuperado de \url{https://www.microsoft.com/es-es/security/business/security-101/what-is-information-security-infosec}

\bibitem{IBM2025Cyber}
IBM. (s.f.). \textit{Cybersecurity}. Recuperado de \url{https://www.ibm.com/es-es/topics/cybersecurity}

\bibitem{AWS2025}
AWS. (s.f.). \textit{What is cybersecurity?} Recuperado de \url{https://aws.amazon.com/es/what-is/cybersecurity/}

\bibitem{Computing2025}
Computing. (s.f.). \textit{Seguridad informática: qué es y qué necesito saber}. Recuperado de \url{https://www.computing.es/seguridad/seguridad-informatica-que-es-y-que-necesito-saber/}

\bibitem{IBM2025Info}
IBM. (s.f.). \textit{Information security}. Recuperado de \url{https://www.ibm.com/mx-es/think/topics/information-security}

\end{thebibliography}



\end{document}
